% =====================================================================
%  abstract.tex
% =====================================================================
\chapter*{Abstract}
\addcontentsline{toc}{chapter}{Abstract}
\thispagestyle{plain}

% Colonna di testo più stretta e centrata (rientro di 1.8 cm per lato)
\begin{list}{}{%
    \setlength{\leftmargin}{1.5cm}%
    \setlength{\rightmargin}{1.5cm}%
    \setlength{\itemsep}{0pt}%
    \setlength{\parsep}{0.8em}%
    \setlength{\topsep}{0pt}}
\item[]
{\color{accent}\rule{\linewidth}{0.8pt}}

\vspace{0.6em}
This document describes the functional architecture of a
\emph{software-as-a-service} application that forecasts the glucose level of
patients with Type~1 diabetes thirty minutes after a meal.

Before eating, the patient declares the carbohydrates of the meal. A machine
learning model, fed with the patient's continuous glucose measurements,
predicts the glucose level, and the patient immediately receives one of four
outcomes: \emph{proceed}, \emph{proceed with a warning}, \emph{risk of
hyperglycemia} or \emph{risk of hypoglycemia}. Risky predictions are recorded
for the doctor, who reviews them periodically together with the trend of each
patient, while the center manager follows aggregated statistics on the use and
the performance of the service.

The document identifies the three actors and the data to acquire, and lists the
functional and non-functional requirements derived from the Service Level
Agreement: service from 6:00 to 22:00, availability and reliability of at least
99.5\%, response time within one second, 50\,000 customers and 5\,000 requests
per minute. It then describes eight use cases and illustrates the behavior of
the system through sequence diagrams. Implementation choices are deliberately
left to the following design phases.

\vspace{0.4em}
{\color{accent}\rule{\linewidth}{0.8pt}}

\vspace{0.2em}


\end{list}